\documentclass[10pt]{article}
\input{cheatsheets/src/cheatsheet-preamble.tex}

\begin{document}
\cheatsheettitle{Introductory Trigonometry Formula Sheet}{Unit Circle, Right Triangles, Periodicity, and Identities}
\small

\begin{minipage}[t]{0.49\textwidth}
\textbf{Unit Circle}\par
At angle $\theta$, the point is
\[
(x,y)=(\cos\theta,\sin\theta),\qquad \tan\theta=\frac{y}{x}.
\]

\centering
\renewcommand{\arraystretch}{1.1}
\begin{tabular}{c c c}
\toprule
Degrees & Radians & $(\cos\theta,\sin\theta)$\\
\midrule
$0^\circ$ & $0$ & $(1,0)$\\
$30^\circ$ & $\frac{\pi}{6}$ & $(\frac{\sqrt3}{2},\frac12)$\\
$45^\circ$ & $\frac{\pi}{4}$ & $(\frac{\sqrt2}{2},\frac{\sqrt2}{2})$\\
$60^\circ$ & $\frac{\pi}{3}$ & $(\frac12,\frac{\sqrt3}{2})$\\
$90^\circ$ & $\frac{\pi}{2}$ & $(0,1)$\\
$120^\circ$ & $\frac{2\pi}{3}$ & $(-\frac12,\frac{\sqrt3}{2})$\\
$135^\circ$ & $\frac{3\pi}{4}$ & $(-\frac{\sqrt2}{2},\frac{\sqrt2}{2})$\\
$150^\circ$ & $\frac{5\pi}{6}$ & $(-\frac{\sqrt3}{2},\frac12)$\\
$180^\circ$ & $\pi$ & $(-1,0)$\\
$210^\circ$ & $\frac{7\pi}{6}$ & $(-\frac{\sqrt3}{2},-\frac12)$\\
$225^\circ$ & $\frac{5\pi}{4}$ & $(-\frac{\sqrt2}{2},-\frac{\sqrt2}{2})$\\
$240^\circ$ & $\frac{4\pi}{3}$ & $(-\frac12,-\frac{\sqrt3}{2})$\\
$270^\circ$ & $\frac{3\pi}{2}$ & $(0,-1)$\\
$300^\circ$ & $\frac{5\pi}{3}$ & $(\frac12,-\frac{\sqrt3}{2})$\\
$315^\circ$ & $\frac{7\pi}{4}$ & $(\frac{\sqrt2}{2},-\frac{\sqrt2}{2})$\\
$330^\circ$ & $\frac{11\pi}{6}$ & $(\frac{\sqrt3}{2},-\frac12)$\\
$360^\circ$ & $2\pi$ & $(1,0)$\\
\bottomrule
\end{tabular}

\raggedright
\textbf{Quadrant signs}\par
QI: all positive; QII: sine; QIII: tangent; QIV: cosine.

\textbf{Angle conversion}
\[
180^\circ=\pi,\qquad
\text{degrees}\cdot\frac{\pi}{180}=\text{radians}.
\]
\end{minipage}
\hfill
\begin{minipage}[t]{0.47\textwidth}
\textbf{SOH--CAH--TOA}
\[
\sin\theta=\frac{\text{opp}}{\text{hyp}},\qquad
\cos\theta=\frac{\text{adj}}{\text{hyp}},\qquad
\tan\theta=\frac{\text{opp}}{\text{adj}}.
\]
\[
\csc\theta=\frac1{\sin\theta},\qquad
\sec\theta=\frac1{\cos\theta},\qquad
\cot\theta=\frac1{\tan\theta}.
\]

\textbf{Pythagorean theorem and identities}
\[
a^2+b^2=c^2
\]
\[
\sin^2\theta+\cos^2\theta=1
\]
\[
1+\tan^2\theta=\sec^2\theta,
\qquad
1+\cot^2\theta=\csc^2\theta.
\]

\textbf{Fundamental periods}
\[
\begin{array}{c|c}
\text{Functions} & \text{Period}\\ \hline
\sin,\ \cos,\ \sec,\ \csc & 2\pi\\
\tan,\ \cot & \pi
\end{array}
\]
If $f(x)$ has period $P$, then $f(bx)$ has period
\[
\frac{P}{|b|}.
\]
Thus $f(x+P)=f(x)$ whenever both sides are defined.

\textbf{Complement (cofunction) rules}\par
For $\theta$ in degrees, replace $\frac\pi2-\theta$ below by $90^\circ-\theta$.
\[
\sin\left(\frac\pi2-\theta\right)=\cos\theta,
\qquad
\cos\left(\frac\pi2-\theta\right)=\sin\theta
\]
\[
\tan\left(\frac\pi2-\theta\right)=\cot\theta,
\qquad
\cot\left(\frac\pi2-\theta\right)=\tan\theta
\]
\[
\sec\left(\frac\pi2-\theta\right)=\csc\theta,
\qquad
\csc\left(\frac\pi2-\theta\right)=\sec\theta.
\]

\textbf{Undefined values}\par
A reciprocal function is undefined when its denominator is zero:
\[
\tan,\sec:\ \cos\theta=0,
\qquad
\cot,\csc:\ \sin\theta=0.
\]
\end{minipage}

\end{document}
